\section{Alcance de la propuesta}

El presente proyecto tiene como alcance el diseño, escalamiento y validación de una tecnología de secado de café basada en el control de humedad relativa mediante bomba de calor, orientada a su aplicación en procesos agroindustriales para la estandarización de cafés diferenciados.

En el marco del programa de Maestría en Ingeniería Agroindustrial de la Universidad Nacional de Colombia, sede Palmira, la propuesta contempla el desarrollo de un sistema de secado a nivel piloto avanzado, con criterios de diseño que permitan su escalabilidad a capacidades industriales comprendidas entre 5 y 20 toneladas por lote. El alcance incluye el modelado del proceso de secado, la optimización de variables operativas, la implementación de estrategias de manejo del producto, así como la evaluación de la calidad del café obtenido bajo diferentes condiciones de operación.

Asimismo, el proyecto considera la validación experimental de la tecnología en condiciones controladas y su proyección hacia escenarios reales de aplicación, particularmente en plantas de beneficio centralizadas como la propuesta en el municipio de Sevilla, Valle del Cauca. En este contexto, se analizará la viabilidad técnica, económica y social de su implementación, así como su integración con modelos productivos como el esquema de compra de café en cereza (CERPER).

En términos de articulación con la Política de Investigación e Innovación Orientada por Misiones (PIIOM), el alcance del proyecto se alinea con la Misión de Bioeconomía y Territorio, al promover la valorización del café mediante el desarrollo de tecnologías que permiten diversificar la oferta de productos y mejorar su competitividad en mercados internacionales. De igual manera, se vincula con la Misión de Energía Sostenible, al incorporar el uso de tecnologías eficientes como la bomba de calor y la posible integración de fuentes de energía renovable.

Adicionalmente, el proyecto contribuye a la Misión de Ciencia para la Paz, al proponer soluciones que fortalecen las economías lícitas en territorios rurales, mediante el mejoramiento de la calidad del producto, la reducción de pérdidas en la cadena productiva y la generación de mayores ingresos para los caficultores.

Es importante señalar que el alcance de la propuesta se limita al desarrollo, validación y evaluación de la tecnología a nivel piloto escalable, así como al análisis de su viabilidad, sin contemplar su implementación a escala industrial completa dentro del tiempo de ejecución de la maestría. No obstante, se espera que los resultados obtenidos sirvan como base para futuros procesos de transferencia tecnológica e implementación a gran escala.